\begin{textsample}{NEG Dim 4 – vem\_ed\_17 – Score 5.00 – vem\_ed\_17/t075.txt}  \label{ex:f4_neg_085}
Fazer COIL permite enxergar que há pessoas no mundo que são como eu, permite ver questões como as consequências da escravidão no Brasil e nos Estados Unidos.
Os estudantes merecem trazer sua cultura para a sala de aula.
Parte do que vamos fazer nos próximos cinco anos envolve o pensamento inclusivo e a interseccionalidade.
Como falamos sobre ações sociais, ambientais e econômicas de pessoas e comunidades, o que fazemos nos negócios afeta o clima e a todos nós.
Outros aspectos envolvem as desigualdades sociais e a maior participação de pessoas que não são brancas e não falam inglês.
Há medo e é preciso coragem, é difícil e então há uma maneira de ter conversas e experiências incríveis e reconhecer que cada um de nós pode estar no centro.
Acho que o COIL vai mudar nos próximos cinco anos se mais pessoas fizerem projetos no ensino fundamental e médio: trabalhar por meio das culturas cada vez mais cedo.
Um dos objetivos que estou buscando é obter informações na comunidade COIL – Centro Paula Souza, SUNY, Florida International University, DePaul, enfim, onde quer que possamos descobrir como fazer crowdsourcing e compartilhar essas informações com toda a comunidade, para que não fiquem isoladas em silos.
Quero que possamos usar big data e extrair informações para os projetos COIL.
Que possamos usar hashtags, wikis ou formas de compartilhar informação por meio das comunidades e aprender com isso.
E isso é algo que eu, como designer instrucional, sei que é possível.
Cronologicamente falando, quais foram os avanços mais significativos dos projetos COIL?
No início dos anos 2000, usávamos o Skype e estávamos apenas tentando descobrir como fazer essa coisa de videoconferência e quais ferramentas on-line usar.
Nessa primeira onda, a pergunta era: como fazer?
Depois, entre 2017, 2018 ou mesmo 2019, a pergunta passou a ser o que vamos fazer?
Quando veio a pandemia, as pessoas começaram a se interessar pelo COIL como uma alternativa para expandir as possibilidades de internacionalização para mais estudantes e incluir as pessoas com menos recursos financeiros, como uma comunidade de prática que poderia ser incluída na sala de aula em vez de uma atividade extracurricular.
Viajar, apenas 1 \% consegue.
A terceira onda, de agora em diante, envolve a questão: por que fazemos COIL?
Como abordamos e incluímos diferenças na nossa comunidade, aprendendo sobre os outros e como superamos essas diferenças?
Com a pandemia e o assassinato de George Floyd, foi preciso repensar, reposicionar todo o programa internacional.
O que os programas internacionais e de educação global fornecem de melhor?
São conversas sobre cultura e aprendizado experiencial.
Portanto, é preciso consolidar a confiança e buscar o"timing"adequado com base na comunicação.
Isso é o que eu amo tanto no COIL: é uma oportunidade entusiasmante, mesmo quando tantas coisas dão errado.

% matched lemmas: 
\end{textsample}
